% GENERATED FILE. Edit canonical lessons, then run npm run generate:books.
% canonical-source-hash: fnv1a64:a165ec830fe1bac1
% canonical-lessons: RU-C109-lawyer, RU-C109-engineer, RU-C109-firefighter, RU-C109-postman, RU-C109-waiter

\chapter{Lawyer, Engineer, Firefighter, Postman, Waiter}
\label{ch:ru-lawyer-engineer-firefighter-postman-waiter}
\hlchaptermodality{109}

\begin{chapteropening}
\textbf{By the end of this chapter:} \emph{I can say a lawyer, an engineer, a firefighter, a postman and a waiter.}

Complete the last lesson of chapter 109: i can say a lawyer, an engineer, a firefighter, a postman and a waiter.
\end{chapteropening}

\section[yuríst]{\ru{юрист} (yuríst) — a lawyer}

\label{lesson:RU-C109-lawyer}



\begin{quote}
Before the new one: say the Russian for a baker, then the Russian for a driver.
\end{quote}

\subsection*{You'll want to know: \ru{юрист}}
\textbf{\ru{юрист}} (\emph{yuríst}) — \textquotedblleft{}a lawyer\textquotedblright{}.

\begin{grammarlens}[title={What you've built}]
One more word for this chapter.
\end{grammarlens}

\subsection*{Your turn}
\begin{itemize}
\raggedright
  \item \emph{Say it:} \emph{yuríst}
  \item \emph{Say it:} \emph{yuríst}, once more
  \item \emph{Say it:} say \emph{yuríst}
  \item \emph{Recall:} say the Russian for a baker, then the Russian for a driver, then say \emph{yuríst} again
\end{itemize}

\begin{tcolorbox}[breakable,colback=teal!4,colframe=teal!35!black,title={Before you move on}]
What does \ru{юрист} mean? (\textquotedblleft{}A lawyer\textquotedblright{}.) Say it once more.
\end{tcolorbox}
\section[inzhenér]{\ru{инженер} (inzhenér) — an engineer}

\label{lesson:RU-C109-engineer}



\begin{quote}
Before the new one: say the Russian for a driver, then the Russian for a lawyer.
\end{quote}

\subsection*{You'll want to know: \ru{инженер}}
\textbf{\ru{инженер}} (\emph{inzhenér}) — \textquotedblleft{}an engineer\textquotedblright{}.

\begin{grammarlens}[title={What you've built}]
One more word for this chapter.
\end{grammarlens}

\subsection*{Your turn}
\begin{itemize}
\raggedright
  \item \emph{Say it:} \emph{inzhenér}
  \item \emph{Say it:} \emph{inzhenér}, once more
  \item \emph{Say it:} say \emph{inzhenér}
  \item \emph{Recall:} say the Russian for a driver, then the Russian for a lawyer, then say \emph{inzhenér} again
\end{itemize}

\begin{tcolorbox}[breakable,colback=teal!4,colframe=teal!35!black,title={Before you move on}]
What does \ru{инженер} mean? (\textquotedblleft{}An engineer\textquotedblright{}.) Say it once more.
\end{tcolorbox}
\section[pozhárnyy]{\ru{пожарный} (pozhárnyy) — a firefighter}

\label{lesson:RU-C109-firefighter}



\begin{quote}
Before the new one: say the Russian for a lawyer, then the Russian for an engineer.
\end{quote}

\subsection*{You'll want to know: \ru{пожарный}}
\textbf{\ru{пожарный}} (\emph{pozhárnyy}) — \textquotedblleft{}a firefighter\textquotedblright{}.

\begin{grammarlens}[title={What you've built}]
One more word for this chapter.
\end{grammarlens}

\subsection*{Your turn}
\begin{itemize}
\raggedright
  \item \emph{Say it:} \emph{pozhárnyy}
  \item \emph{Say it:} \emph{pozhárnyy}, once more
  \item \emph{Say it:} say \emph{pozhárnyy}
  \item \emph{Recall:} say the Russian for a lawyer, then the Russian for an engineer, then say \emph{pozhárnyy} again
\end{itemize}

\begin{tcolorbox}[breakable,colback=teal!4,colframe=teal!35!black,title={Before you move on}]
What does \ru{пожарный} mean? (\textquotedblleft{}A firefighter\textquotedblright{}.) Say it once more.
\end{tcolorbox}
\section[pochtalón]{\ru{почтальон} (pochtalón) — a postman}

\label{lesson:RU-C109-postman}



\begin{quote}
Before the new one: say the Russian for an engineer, then the Russian for a firefighter.
\end{quote}

\subsection*{You'll want to know: \ru{почтальон}}
\textbf{\ru{почтальон}} (\emph{pochtalón}) — \textquotedblleft{}a postman\textquotedblright{}.

\begin{grammarlens}[title={What you've built}]
One more word for this chapter.
\end{grammarlens}

\subsection*{Your turn}
\begin{itemize}
\raggedright
  \item \emph{Say it:} \emph{pochtalón}
  \item \emph{Say it:} \emph{pochtalón}, once more
  \item \emph{Say it:} say \emph{pochtalón}
  \item \emph{Recall:} say the Russian for an engineer, then the Russian for a firefighter, then say \emph{pochtalón} again
\end{itemize}

\begin{tcolorbox}[breakable,colback=teal!4,colframe=teal!35!black,title={Before you move on}]
What does \ru{почтальон} mean? (\textquotedblleft{}A postman\textquotedblright{}.) Say it once more.
\end{tcolorbox}
\section[ofitsiánt]{\ru{официант} (ofitsiánt) — a waiter}

\label{lesson:RU-C109-waiter}



\begin{quote}
Before the new one: say the Russian for a firefighter, then the Russian for a postman.
\end{quote}

\subsection*{You'll want to know: \ru{официант}}
\textbf{\ru{официант}} (\emph{ofitsiánt}) — \textquotedblleft{}a waiter\textquotedblright{}.

\begin{grammarlens}[title={What you've built}]
One more word for this chapter.
\end{grammarlens}

\subsection*{Your turn}
\begin{itemize}
\raggedright
  \item \emph{Say it:} \emph{ofitsiánt}
  \item \emph{Say it:} \emph{ofitsiánt}, once more
  \item \emph{Say it:} say \emph{ofitsiánt}
  \item \emph{Recall:} say the Russian for a firefighter, then the Russian for a postman, then say \emph{ofitsiánt} again
\end{itemize}

\begin{tcolorbox}[breakable,colback=teal!4,colframe=teal!35!black,title={Before you move on}]
What does \ru{официант} mean? (\textquotedblleft{}A waiter\textquotedblright{}.) Say it once more.
\end{tcolorbox}
